\documentclass[12pt,a4paper]{report}

% ============================================================
% PACKAGES
% ============================================================
\usepackage[a4paper,left=1.5in,right=1in,top=1in,bottom=1in]{geometry}
\usepackage{times}
\usepackage[T1]{fontenc}
\usepackage{graphicx}
\usepackage{fancyhdr}
\usepackage{titlesec}
\usepackage{tocloft}
\usepackage{setspace}
\usepackage{xcolor}
\usepackage{longtable}
\usepackage{array}
\usepackage{booktabs}
\usepackage{enumitem}
\usepackage{listings}
\usepackage{amsmath}
\usepackage{hyperref}
\usepackage{caption}
\usepackage{float}

\definecolor{navy}{HTML}{0F172A}
\definecolor{brand}{HTML}{0EA5E9}
\definecolor{codebg}{HTML}{F8FAFC}

\hypersetup{
  colorlinks=true,
  linkcolor=navy,
  citecolor=navy,
  urlcolor=brand,
  pdftitle={TransitSwap: A Transparent, Accessibility-Aware Multimodal Transit Route Planner},
  pdfauthor={Srujan V, AL Abhishek, Shreyas M, B Anand}
}

% ---- Chapter / section heading style (centered bold caps, matching the reference) ----
\titleformat{\chapter}[display]
  {\centering\normalfont\bfseries\Large}
  {\chaptertitlename\ \thechapter}{8pt}{\Large\uppercase}
\titlespacing*{\chapter}{0pt}{-20pt}{30pt}
\titleformat{\section}{\normalfont\bfseries\large}{\thesection}{8pt}{}
\titleformat{\subsection}{\normalfont\bfseries\normalsize}{\thesubsection}{6pt}{}

% ---- Header / footer, matching the reference's running-title + page-number style ----
\setlength{\headheight}{15pt}
\pagestyle{fancy}
\fancyhf{}
\renewcommand{\headrulewidth}{1.2pt}
\renewcommand{\footrulewidth}{0.6pt}
\fancyhead[L]{\small\bfseries TransitSwap: A Transparent, Accessibility-Aware Multimodal Transit Route Planner}
\fancyfoot[C]{\small\thepage}
\fancyfoot[L]{\small 2025--26}

\setstretch{1.25}
\setlength{\cftbeforechapskip}{8pt}
\renewcommand{\cftchapfont}{\bfseries}

% ---- Listings style for pseudocode / real code excerpts ----
\lstset{
  basicstyle=\ttfamily\footnotesize,
  backgroundcolor=\color{codebg},
  frame=single,
  rulecolor=\color{gray!40},
  breaklines=true,
  showstringspaces=false,
  columns=flexible,
  numbers=left,
  numberstyle=\tiny\color{gray},
  xleftmargin=1.8em,
}

% ---- Placeholder figure box, used throughout Chapter 7 where the author will
%      paste in real application screenshots before submission ----
\newcommand{\screenshotplaceholder}[2]{%
  \begin{figure}[H]
    \centering
    \fbox{\begin{minipage}[c][6cm][c]{0.85\textwidth}
      \centering
      \textit{\textcolor{gray}{[ Insert screenshot here: #2 ]}}
    \end{minipage}}
    \caption{#1}
  \end{figure}
}

\begin{document}

% ============================================================
% TITLE PAGE
% ============================================================
\begin{titlepage}
\begin{center}

{\Large \textbf{<YOUR UNIVERSITY NAME HERE>}}\\[2pt]
{\small ``<University motto / address>''}\\[30pt]

{\large \textbf{PROJECT REPORT}}\\[4pt]
On\\[10pt]

{\Huge \textbf{TransitSwap}}\\[6pt]
{\Large \textbf{A Transparent, Accessibility-Aware Multimodal}}\\
{\Large \textbf{Transit Route Planner}}\\[20pt]

{\normalsize Submitted in partial fulfillment of the requirements of the degree of}\\[6pt]
{\large \textbf{BACHELOR OF ENGINEERING}}\\[4pt]
in\\[4pt]
{\large \textbf{Information Science and Engineering}}\\[6pt]
{\normalsize For the academic year}\\
{\normalsize 2025--2026}\\[24pt]

{\normalsize Submitted By}\\[6pt]
{\large \textbf{Srujan V}} \quad (\textbf{<USN>})\\
{\large \textbf{AL Abhishek}} \quad (\textbf{<USN>})\\
{\large \textbf{Shreyas M}} \quad (\textbf{<USN>})\\
{\large \textbf{B Anand}} \quad (\textbf{<USN>})\\[20pt]

{\normalsize Under the guidance of}\\[6pt]
{\large \textbf{Swathi B Patil}}\\
{\normalsize Assistant Professor, Dept. of Information Science and Engineering,}\\
{\normalsize Don Bosco Institute of Technology}\\[30pt]

{\large \textbf{DEPARTMENT OF INFORMATION SCIENCE AND ENGINEERING}}\\[4pt]
{\large \textbf{DON BOSCO INSTITUTE OF TECHNOLOGY}}\\[2pt]
{\normalsize <Affiliation line, e.g. An Autonomous Institute, Affiliated to University>}\\
{\normalsize Bengaluru, Karnataka, India}\\

\end{center}
\end{titlepage}

% ============================================================
% CERTIFICATE
% ============================================================
\thispagestyle{empty}
\begin{center}
{\large \textbf{DON BOSCO INSTITUTE OF TECHNOLOGY}}\\
{\small <Affiliation line>}\\
{\small Bengaluru, Karnataka, India}\\[16pt]
{\large \textbf{DEPARTMENT OF INFORMATION SCIENCE AND ENGINEERING}}\\[12pt]
{\Large \textbf{CERTIFICATE}}
\end{center}
\vspace{10pt}
\noindent
Certified that the project work entitled \textbf{``TransitSwap: A Transparent, Accessibility-Aware
Multimodal Transit Route Planner''} is a bonafide work carried out by \textbf{Srujan V, AL Abhishek, Shreyas M and B Anand},
bearing \textbf{USNs <USN>, <USN>, <USN> and <USN>} respectively, bonafide students of \textbf{Don Bosco Institute of Technology}, in partial
fulfillment for the award of the degree of Bachelor of Engineering in \textbf{Information Science and Engineering} of
\textbf{<University Name>} during the academic year 2025--26. It is certified that all
corrections/suggestions indicated during the internal assessment have been incorporated in the
report deposited in the departmental library. The project report has been approved as it
satisfies the academic requirements with respect to the project work prescribed for the said
degree.

\vspace{40pt}
\noindent
\begin{minipage}{0.3\textwidth}
\centering
Signature of the Guide\\[20pt]
\textbf{Swathi B Patil}\\
Assistant Professor
\end{minipage}
\hfill
\begin{minipage}{0.3\textwidth}
\centering
Signature of the HOD\\[20pt]
\textbf{<HOD Name>}\\
Professor \& Head
\end{minipage}
\hfill
\begin{minipage}{0.3\textwidth}
\centering
Signature of the Principal\\[20pt]
\textbf{<Principal Name>}\\
Principal
\end{minipage}

\vspace{30pt}
\noindent
\textbf{Name of the Examiners} \hfill \textbf{Signature with Date}\\[10pt]
1. \underline{\hspace{12cm}}\\[8pt]
2. \underline{\hspace{12cm}}

\clearpage

% ============================================================
% DECLARATION
% ============================================================
\thispagestyle{empty}
\begin{center}
{\Large \textbf{DECLARATION}}
\end{center}
\vspace{10pt}
\noindent
We, \textbf{Srujan V, AL Abhishek, Shreyas M and B Anand}, final-year students of Bachelor of Engineering
in \textbf{Information Science and Engineering}, \textbf{Don Bosco Institute of Technology}, hereby declare that the project work
entitled \textbf{``TransitSwap: A Transparent, Accessibility-Aware Multimodal Transit Route
Planner''} submitted to \textbf{<University Name>} during the academic year 2025--26 is a record
of an original work done by us under the guidance of \textbf{Swathi B Patil}, \textbf{Assistant Professor},
Department of \textbf{Information Science and Engineering}, \textbf{Don Bosco Institute of Technology}. This project work is submitted in
partial fulfilment of the requirements for the award of the degree of Bachelor of Engineering in
\textbf{Information Science and Engineering}. The results embodied in this project have not been submitted to any other
University or Institute for the award of any degree.

\vspace{30pt}
\noindent
\begin{tabular}{@{}lll@{}}
\textbf{Name} & \textbf{USN} & \textbf{Signature} \\[8pt]
Srujan V & <USN> & \underline{\hspace{4cm}} \\[6pt]
AL Abhishek & <USN> & \underline{\hspace{4cm}} \\[6pt]
Shreyas M & <USN> & \underline{\hspace{4cm}} \\[6pt]
B Anand & <USN> & \underline{\hspace{4cm}} \\
\end{tabular}

\clearpage

% ============================================================
% ACKNOWLEDGEMENT
% ============================================================
\thispagestyle{empty}
\begin{center}
{\Large \textbf{ACKNOWLEDGEMENT}}
\end{center}
\vspace{10pt}
\noindent
The satisfaction and euphoria that accompany the successful completion of any task would be
incomplete without the mention of the people who made it possible, whose constant guidance and
encouragement crowned this effort with success.

\vspace{8pt}\noindent
We are grateful to our institution, \textbf{Don Bosco Institute of Technology}, for its ideals and inspiration, and for
providing us with the facilities that made this project possible.

\vspace{8pt}\noindent
We sincerely thank \textbf{<Principal Name>}, Principal, \textbf{Don Bosco Institute of Technology}, for facilitating
academic excellence and providing a congenial environment to work in.

\vspace{8pt}\noindent
We wish to extend our profound thanks to \textbf{<HOD Name>}, Professor \& Head, Department of
\textbf{Information Science and Engineering}, for granting consent to carry out this project.

\vspace{8pt}\noindent
We express our sincere thanks to our guide, \textbf{Swathi B Patil}, \textbf{Assistant Professor}, Department
of \textbf{Information Science and Engineering}, for the able guidance and valuable advice at every stage, which helped in
the successful completion of this project.

\vspace{8pt}\noindent
I would like to thank all the teaching and non-teaching staff for their valuable advice and
support, and we express our sincere thanks to our parents and friends for their continued
encouragement.

\vspace{30pt}
\noindent Srujan V, AL Abhishek, Shreyas M, B Anand

\clearpage

% ============================================================
% ABSTRACT
% ============================================================
\thispagestyle{empty}
\begin{center}
{\Large \textbf{ABSTRACT}}
\end{center}
\vspace{10pt}
\noindent
Mainstream navigation applications optimise a route almost entirely for travel time and expose
little else to the commuter, and none of the major consumer tools let a rider express a genuine
hard constraint such as ``this route must be step-free end to end'' -- accessibility, if offered at
all, is only ever a soft preference that a faster option can silently override. \textbf{TransitSwap}
is a prototype multimodal journey planner built to close that gap. It ranks every candidate journey
on nine normalised features -- time, cost, walking distance, reliability, accessibility, crowding,
weather, missed-connection risk and transfer count -- using a linear scorer whose weights are
learned per user from recorded pairwise route choices via a pairwise logistic regression, with
profile-specific presets for seven mobility profiles (standard, wheelchair, senior, pregnant,
stroller, luggage, reduced mobility). Accessibility is enforced as a genuine \emph{hard} filter per
profile, not merely a weighted score, and every recommendation is accompanied by a plain-language
explanation whose every claim is checked against the real candidate set before being shown, rather
than asserted. The system is implemented end to end on a React~19~/~TypeScript frontend and an
Express~/~TypeScript~/~MongoDB backend: a bounded round-based graph search generates candidate
journeys over a seeded Bengaluru transit network; real road geometry is obtained from OSRM and
real place search from Nominatim; weather, crowd and historical-delay signals feed a rule-based
reliability score and a data-driven missed-connection-risk Monte~Carlo simulation; and a
50-endpoint REST API is exercised by 29 backend and 35 frontend automated tests, all passing.
The resulting system was evaluated on a 22-scenario benchmark against five baselines, and an
ablation study confirms every non-time feature has real, measurable influence on the final ranking.
This report documents the complete design, implementation, and evaluation of TransitSwap --
including its honestly reported limitations, since the transit network itself is a synthetic,
hand-authored dataset and no production-scale user study was conducted -- as a final-year major
project.

\clearpage

% ============================================================
% TOC / LOF / LOT
% ============================================================
\pagenumbering{roman}
\tableofcontents
\clearpage
\listoffigures
\clearpage
\listoftables
\clearpage

\pagenumbering{arabic}
\setcounter{page}{1}

% ============================================================
% CHAPTER 1 — INTRODUCTION
% ============================================================
\chapter{Introduction}

\section{Briefing About the Project}

Public transport in a large Indian city is rarely a single mode. A real journey typically combines
a short walk, a metro ride, possibly a bus leg, and sometimes a last-mile auto-rickshaw -- and the
right combination for one commuter is frequently the wrong one for another. A wheelchair user
cannot take a route through a station with a broken lift and uncovered stairs. A commuter travelling
with a stroller or heavy luggage experiences a crowded platform very differently from an
unencumbered rider. A senior citizen's tolerance for a long walk or a tight interchange is not the
same as a student's. Mainstream navigation tools largely ignore this: they compute the fastest or
shortest route and, at best, offer accessibility as an afterthought.

TransitSwap is a prototype multimodal journey planner designed around the opposite premise: that a
transparent, multi-criteria ranking -- with accessibility enforced as a real constraint, not a
suggestion, and every recommendation explained in plain language -- is more useful to a commuter
than a single time-optimal answer, \emph{provided} the ranking genuinely produces different
decisions for different people and every claim the system makes is checked against real data
before being shown. The project implements this premise completely: real road-snapped walking and
auto-rickshaw geometry (via OSRM), a real round-based graph search over a seeded Bengaluru transit
network for metro and bus candidates, a genuine per-user learned ranking model (a pairwise logistic
regression trained on the user's own recorded route choices), a hard accessibility-constraint
engine covering seven distinct mobility profiles, and a reliability layer that is honest about what
it does and does not know -- refusing to report a confidence level, for example, when too little
historical data exists to compute one legitimately.

\section{Existing System}

The dominant consumer routing applications in use today -- Google Maps, Citymapper, and similar
tools -- are built around a single-criterion or lightly-weighted multi-criterion shortest-path
engine, almost always optimised primarily for travel time.

\subsection{Use of Static Image-Based Systems}
Several research and campus-level transit tools still rely on static, pre-computed timetable
lookups rather than a live graph search, which cannot adapt to a live network state, a closed
station, or a user-specific constraint at query time.

\subsection{Use of Hardware-Based Systems}
Some accessibility-focused transit aids depend on dedicated hardware -- beacons, RFID tags at
station entrances, or proprietary sensor networks -- to determine real-time accessibility status.
This raises deployment cost and ties the system to a specific transit authority's hardware
investment, which most mid-sized city networks do not have.

\subsection{Vocabulary and Training Data}
Existing accessible-routing research prototypes are frequently trained or tuned on a narrow,
single-city dataset and do not generalise their constraint logic across multiple rider profiles;
a system designed for wheelchair users alone, for instance, typically has no equivalent logic for
a pregnant commuter or a rider travelling with a stroller.

\subsection{Real-Time Processing}
Most existing multi-criteria research prototypes are evaluated offline, against a fixed scenario
set, and are not wired into a live, queryable API that a real frontend application can call under
interactive latency constraints.

\section{Limitations of Existing System}

\begin{itemize}[leftmargin=*]
  \item Accessibility, where offered, is a soft preference that a faster or cheaper route can
    silently override -- there is no hard guarantee that a wheelchair-profile route is actually
    step-free throughout.
  \item Crowd level and walking distance, even where estimated, rarely carry enough ranking weight
    to change the final recommendation, so a commuter with a real physical constraint still sees
    the same top route as everyone else.
  \item No plain-language explanation is given for why one route was preferred over another, so the
    ranking is effectively a black box to the end user.
  \item Personalisation, where present, is usually all-or-nothing: once a system has learned a
    user's general preferences, a profile selected for that specific trip (e.g. ``travelling with a
    stroller today'') can be silently ignored.
  \item Reliability and delay estimates are frequently presented with a false sense of precision --
    a stated confidence level is shown even when the underlying historical sample is too small to
    support it honestly.
\end{itemize}

\section{Proposed System}

TransitSwap addresses each of the limitations above directly.

\begin{itemize}[leftmargin=*]
  \item \textbf{A genuine hard accessibility filter.} Seven mobility profiles each carry their own
    rule set over real per-station accessibility attributes (lift, ramp, step-free entrance, stair
    count); a route that fails a profile's hard constraint is removed from the candidate set
    entirely, not merely down-ranked.
  \item \textbf{A learned, per-user multi-criteria ranker.} Each candidate route is scored as a
    weighted linear sum over nine normalised features; the weights start from a profile-specific
    preset and are refined per user via a pairwise logistic regression trained on the user's own
    recorded choices, with a floor that guarantees a profile's safety-relevant weighting (crowd,
    walking, accessibility) is never diluted below its own preset value, even once personalisation
    has taken over.
  \item \textbf{An honesty-gated explanation layer.} Every recommendation carries plain-language
    tags (``fastest option'', ``lowest estimated fare'', ``minimizes walking distance'') -- each one
    individually checked against the real candidate set before being shown, so a claim is never
    displayed unless it is actually true for that route.
  \item \textbf{Statistically grounded, honestly gated reliability.} A data-driven prediction
    interval and a Monte~Carlo missed-connection-risk simulation are both computed from real
    historical observation data where it exists, and both explicitly refuse to report a confidence
    level when too little data exists to support one.
  \item \textbf{Real, live external data.} Road-snapped walking and driving geometry comes from a
    live OSRM instance; place search and reverse geocoding come from Nominatim; current weather
    comes from OpenWeatherMap where configured. None of this is invented or mocked.
\end{itemize}

\begin{figure}[H]
  \centering
  \includegraphics[width=0.95\textwidth]{figures/fig-4.0-architecture.png}
  \caption{Proposed system -- layered architecture overview}
\end{figure}

\section{Objectives}

\begin{enumerate}[leftmargin=*]
  \item To design and implement a bounded round-based graph search that generates genuinely
    non-dominated multimodal route candidates (walk, metro, bus, auto-rickshaw) over a seeded
    transit network, replacing an earlier, more limited fixed-pattern candidate generator.
  \item To implement a hard accessibility-constraint engine covering seven distinct mobility
    profiles, each with its own rule set over real per-station accessibility attributes.
  \item To implement a multi-criteria ranking model -- a weighted linear sum over nine normalised
    route features -- whose weights are learnable per user via a pairwise logistic regression
    trained on the user's own recorded route choices.
  \item To implement a reliability layer (a rule-based reliability score, a data-driven prediction
    interval, and a Monte~Carlo missed-connection-risk simulation) that is explicit about which of
    its outputs are genuinely data-driven and which are rule-of-thumb fallbacks.
  \item To implement a plain-language explanation layer whose every claim is individually verified
    against the real candidate set before being shown to the user.
  \item To build a complete, production-shaped full-stack application (React/TypeScript frontend,
    Express/TypeScript/MongoDB backend, a 50-endpoint REST API) and to verify its correctness with
    an extensive automated test suite rather than manual spot-checking alone.
  \item To evaluate the system honestly against baseline routing strategies and to report its real
    limitations -- including the synthetic nature of its transit dataset -- rather than overstating
    its readiness for production deployment.
\end{enumerate}

\section{Motivation}

The motivation for this project is twofold. First, accessibility in consumer transit routing is
still treated, in practice, as a secondary feature bolted onto a time-optimal engine, rather than a
first-class constraint -- and the consequence falls disproportionately on commuters who already
face the most friction in using public transport: wheelchair users, pregnant commuters, parents
travelling with small children, and senior citizens. A system that genuinely enforces accessibility
as a hard constraint, and is transparent about why it recommends what it recommends, is a
meaningfully better experience for exactly the commuters who need it most. Second, the project is
motivated by a stricter engineering standard than is typical of student projects: every claim the
system makes to a user -- a confidence interval, a crowd estimate, a ``fastest route'' badge -- is
checked against real, computed data before being shown, and every number reported in this project's
own evaluation traces back to a command that reproduces it, rather than being hand-estimated. This
standard shaped nearly every design decision described in the chapters that follow.

% ============================================================
% CHAPTER 2 — LITERATURE SURVEY
% ============================================================
\chapter{Literature Survey}

\section{System Study}

Before designing TransitSwap, the problem space was studied along three independent threads that
together define the design space for a modern multi-criteria transit router: (1) the algorithmic
literature on multi-criteria and time-table-based transit routing, which defines what an
efficient, correct candidate-generation engine looks like; (2) the learning-to-rank literature,
which defines how a ranking function can be learned from observed user preference rather than
hand-tuned; and (3) the literature on crowd-sourced geographic data and open routing infrastructure,
which defines what realistic, deployable data sources a student project can responsibly build on
without requiring a commercial data licence. Each of these threads directly shaped a concrete design
decision in TransitSwap, documented alongside the relevant citation below.

\section{Review of Literature}

\begin{enumerate}[leftmargin=*]

\item \textbf{Delling, D., Pajor, T., and Werneck, R. F. (2015). ``Round-Based Public Transit
Routing.'' \textit{INFORMS Journal on Computing}, 27(2).}\\
This paper introduces RAPTOR (Round-based Public Transit Optimized Router), the standard algorithm
for computing Pareto-optimal journeys (arrival time, number of transfers) over a full time-table at
city scale, processing one ``round'' per additional transfer rather than running a conventional
Dijkstra-style search. \textbf{Relevance:} TransitSwap's own candidate-generation engine (Chapter~4)
converges on the same underlying idea RAPTOR formalises -- a bounded round-based graph search over a
stop graph that returns a \emph{set} of Pareto-optimal candidates rather than a single shortest
path -- though TransitSwap does not implement full RAPTOR, since it has no real time-table to scan,
only a seeded, headway-approximate network.

\item \textbf{Dibbelt, J., Pajor, T., Strasser, B., and Wagner, D. (2018). ``Connection Scan
Algorithm.'' \textit{ACM Journal of Experimental Algorithmics}, 23.}\\
The Connection Scan Algorithm (CSA) computes earliest-arrival journeys by scanning a sorted list of
individual vehicle connections in a single linear pass, and remains one of the two standard
reference approaches (alongside RAPTOR) for time-table routing at scale. \textbf{Relevance:}
reviewed as the direct algorithmic alternative to the round-based approach TransitSwap's design
draws from; it is specifically named, rather than silently omitted, in this project's own honest
account of what its candidate generator does and does not implement (Chapter~8).

\item \textbf{Burges, C., Shaked, T., Renshaw, E., et al. (2005). ``Learning to Rank using
Gradient Descent.'' \textit{Proceedings of the 22nd International Conference on Machine
Learning (ICML)}.}\\
This paper introduces RankNet, the foundational pairwise learning-to-rank formulation: given a pair
of items and an observed preference between them, learn a scoring function such that the preferred
item scores higher. \textbf{Relevance:} TransitSwap's personalisation model (Section~5.4) is a
direct, lighter-weight instance of this idea -- a linear pairwise logistic regression trained on the
delta between a chosen and a rejected route's feature vectors -- chosen deliberately over a deeper
model because the per-user training-sample volume a route planner can realistically collect (tens,
not millions, of labelled pairs) cannot support anything larger without over-fitting.

\item \textbf{Bast, H., Delling, D., Goldberg, A., et al. (2016). ``Route Planning in
Transportation Networks.'' In \textit{Algorithm Engineering}, Lecture Notes in Computer
Science, Springer.}\\
A comprehensive survey of speed-up techniques and multi-criteria extensions for shortest-path and
transit routing, covering contraction hierarchies, multi-level Dijkstra, and the trade-offs between
query-time latency and preprocessing cost. \textbf{Relevance:} informed the decision to bound
TransitSwap's own graph search (a small \(k\) candidates per query, over a deliberately compact
seeded network) rather than attempt a full preprocessing pipeline that would be disproportionate to
the project's dataset scale.

\item \textbf{Luxen, D., and Vetter, C. (2011). ``Real-time Routing with OpenStreetMap Data.''
\textit{Proceedings of the 19th ACM SIGSPATIAL International Conference on Advances in
Geographic Information Systems (GIS)}.}\\
This paper introduces OSRM (Open Source Routing Machine), the contraction-hierarchy-based routing
engine that computes real road-network shortest paths from OpenStreetMap data at interactive
latency. \textbf{Relevance:} OSRM is used directly, live, as TransitSwap's walking and
auto-rickshaw routing provider (Section~5.3) -- every walking and auto segment's geometry, distance
and duration in this project is real OSRM output, not a straight-line approximation, with a
documented straight-line fallback used only on a live OSRM failure.

\item \textbf{Haklay, M., and Weber, P. (2008). ``OpenStreetMap: User-Generated Street Maps.''
\textit{IEEE Pervasive Computing}, 7(4).}\\
An early, influential study of OpenStreetMap as a crowd-sourced geographic dataset, examining its
coverage, accuracy and the sustainability of volunteer-maintained map data relative to commercial
alternatives. \textbf{Relevance:} underpins the project's choice to build entirely on
OpenStreetMap-derived infrastructure (OSRM for routing, Nominatim for geocoding, and OSM-rendered
basemap tiles for the map UI) -- a choice that keeps the whole system free to run and reproduce,
at the cost of coverage and accuracy limitations the project states honestly rather than conceals.

\item \textbf{Koren, Y., Bell, R., and Volinsky, C. (2009). ``Matrix Factorization Techniques
for Recommender Systems.'' \textit{IEEE Computer}, 42(8).}\\
A widely cited survey of latent-factor and matrix-factorisation approaches to personalised
recommendation, including the cold-start problem -- how a recommender should behave for a user with
no recorded history yet. \textbf{Relevance:} directly informed TransitSwap's own cold-start design
(Section~5.4): a brand-new user is served a sensible profile-specific preset rather than an
undefined or randomly initialised model, and the system only switches to a fully personalised
learned model once a minimum number of real recorded choices exists.

\item \textbf{Guo, Z., and Wilson, N. H. M. (2011). ``Assessing the Cost of Transfer
Inconvenience in Public Transport Systems: A Case Study of the London Underground.''
\textit{Transportation Research Part A}, 45(2).}\\
A transportation-economics study quantifying how much a transfer's perceived cost exceeds its
measured time cost -- riders consistently rate an interchange as more burdensome than the minutes it
actually costs them, and this effect varies meaningfully by passenger type. \textbf{Relevance:}
motivates TransitSwap's explicit \texttt{transfers} feature as a distinct ranking dimension (not
folded into \texttt{time}) and its heavier transfer weighting for the senior, pregnant, luggage and
reduced-mobility profiles specifically, where transfer burden is disproportionately high.

\end{enumerate}

\section{Summary}

The literature reviewed establishes three things the project's own design directly reflects. First,
the standard algorithmic approaches to multi-criteria transit routing (RAPTOR, CSA) are built around
returning a \emph{set} of Pareto-optimal journeys rather than a single shortest path -- exactly the
property TransitSwap's own round-based graph search (Chapter~4) is built to provide, since a
multi-criteria ranker has nothing meaningful to rank over a degenerate, single-candidate set.
Second, the learning-to-rank literature supports a pairwise, linear approach specifically under the
low-data regime a real per-user route planner actually operates in, rather than a deeper model that
would overfit. Third, building entirely on open, crowd-sourced infrastructure (OpenStreetMap, OSRM,
Nominatim) is both a well-studied and a practically sound choice for a project of this scale,
provided -- as this project does throughout -- its coverage and accuracy limitations are stated
honestly rather than implied away.

% ============================================================
% CHAPTER 3 — SYSTEM REQUIREMENTS AND SPECIFICATION
% ============================================================
\chapter{System Requirements and Specification}

\section{Hardware Requirements}

\begin{table}[H]
\centering
\caption{Minimum development/runtime hardware requirements}
\begin{tabular}{@{}ll@{}}
\toprule
\textbf{Component} & \textbf{Requirement} \\
\midrule
Processor & Dual-core 2.0\,GHz or higher (Intel/AMD/Apple Silicon) \\
RAM & 8\,GB minimum (16\,GB recommended for running frontend, backend and \\
    & an in-memory MongoDB instance simultaneously during development) \\
Storage & 2\,GB free disk space (source, dependencies, build artefacts) \\
Network & Broadband internet connection (required for live OSRM, Nominatim, \\
        & OpenWeatherMap, and MongoDB Atlas calls) \\
Display & 1280$\times$720 or higher, for the map-based user interface \\
\bottomrule
\end{tabular}
\end{table}

TransitSwap has no dependency on specialised hardware -- no GPU, no sensor, no proprietary
transit-authority hardware integration -- a deliberate design choice so the system can be deployed
on any standard consumer or cloud-hosted machine.

\section{Software Requirements}

\begin{table}[H]
\centering
\caption{Software stack}
\begin{tabular}{@{}lll@{}}
\toprule
\textbf{Layer} & \textbf{Technology} & \textbf{Version} \\
\midrule
Frontend framework & React & 19.0 \\
Frontend build tool & Vite & 8.0 \\
Frontend language & TypeScript & 5.7 \\
Styling & Tailwind CSS & 4.0 \\
Mapping & Leaflet / react-leaflet & 1.9 / 5.0 \\
Backend runtime & Node.js & 20.x \\
Backend framework & Express & 4.19 \\
Backend language & TypeScript & 5.5 \\
Database & MongoDB (via Mongoose) & 8.4 (driver) \\
Authentication & JSON Web Tokens (jsonwebtoken) & 9.0 \\
Password hashing & bcryptjs & 2.4 \\
HTTP client & Axios & 1.19 \\
Routing engine (external) & OSRM (Open Source Routing Machine) & live public instance \\
Geocoding (external) & Nominatim (OpenStreetMap) & live public instance \\
Weather (external, optional) & OpenWeatherMap API & -- \\
Testing (backend) & Node's built-in test runner + Supertest & -- \\
Testing (frontend) & Vitest + Testing Library & 4.1 \\
End-to-end testing & Playwright & 1.63 \\
Version control & Git / GitHub & -- \\
\bottomrule
\end{tabular}
\end{table}

\section{Functional Requirements}

\begin{enumerate}[leftmargin=*]
  \item The system shall accept an origin and destination (by text search or map interaction) and
    return a set of ranked multimodal route candidates.
  \item The system shall allow the user to select one of seven accessibility profiles, and shall
    exclude any candidate route that fails that profile's hard accessibility constraint.
  \item The system shall rank the remaining candidates using a multi-criteria weighted score, and
    shall learn per-user weights from the user's own recorded route choices once sufficient history
    exists.
  \item The system shall attach a plain-language explanation to each recommended route, with every
    explanation claim individually verified against that route's actual computed data.
  \item The system shall report a data-driven arrival-time confidence interval and a missed
    connection risk, and shall explicitly indicate when either is a rule-of-thumb fallback rather
    than a genuinely data-driven estimate.
  \item The system shall allow a registered user to save a journey, a reusable route, and a named
    destination, and to view their journey history and dashboard statistics.
  \item The system shall allow any user to submit a live crowd-level report or an accessibility
    issue report for a station, and shall allow an administrator to review and act on accessibility
    reports and station records.
  \item The system shall display the selected route on an interactive map with real road-following
    geometry for walking and auto-rickshaw legs, and real station/stop-sequence geometry for metro
    and bus legs.
  \item The system shall degrade gracefully when the database is not configured, continuing to
    serve route generation and accessibility information from seeded demo data rather than failing
    outright.
\end{enumerate}

% ============================================================
% CHAPTER 4 — SYSTEM DESIGN
% ============================================================
\chapter{System Design}

\section{System Architecture}

TransitSwap follows a conventional layered architecture: a React frontend communicates over a REST
API with an Express backend, whose logic is organised strictly as
\texttt{routes~$\rightarrow$~controllers (thin)~$\rightarrow$~services (all domain logic)~$\rightarrow$~models~/~data}.
Every API response uses a single \texttt{\{success, data, message\}} envelope, and every
database-touching service checks the live MongoDB connection state before querying, falling back to
seeded demo data rather than failing when no database is configured (Figure~\ref{fig:arch}).

\begin{figure}[H]
  \centering
  \includegraphics[width=0.95\textwidth]{figures/fig-4.0-architecture.png}
  \caption{System architecture}
  \label{fig:arch}
\end{figure}

\section{Input Module}

The input module is responsible for collecting a journey request from the user: an origin and
destination (resolved either via a text search against the geocoding service, a direct map click, or
the browser's own geolocation API), an optional departure date/time, an accessibility profile
selection, and an optional set of explicit travel preferences (preferred mode, walking tolerance,
budget preference, priority). This is implemented in the frontend's \texttt{PlanTripPage} and the
reusable \texttt{LocationSearch} and \texttt{MapView} components.

\section{Processing Module}

The processing module is the backend's \texttt{multimodalController} and \texttt{multimodalService}:
it validates the incoming request, resolves inactive/deactivated stations, runs candidate generation
(Section~4.7), applies the accessibility hard filter, enriches each surviving candidate with
weather, crowd and reliability context, computes a ranking score for each candidate, and attaches an
explanation to the top-ranked routes before returning the ranked set to the frontend.

\section{Feature Extraction}

Each candidate route is converted into a nine-dimensional normalised feature vector -- time, cost,
walking distance, reliability, accessibility, crowd, weather, missed-connection risk, and transfer
count -- via relative min-max normalisation across the candidate set returned for that query. This
feature vector is the single shared representation consumed by both the rule-based scoring path and
the learned pairwise logistic regression model, implemented in
\texttt{backend/src/services/ml/featureExtractor.ts}.

\section{Classification Module}

The ``classification'' step in TransitSwap is a ranking, not a discrete classification: each
candidate's feature vector is scored as a weighted linear sum, with weights drawn from a
profile-specific preset and refined, per user, by a pairwise logistic regression trained on that
user's own recorded choices (\texttt{mlPreferenceService.ts}). The highest-scoring candidate is
surfaced as the primary recommendation; the full ranked set is also returned so the user can inspect
and select an alternative.

\section{Output Module}

The output module renders the ranked route set as an interactive map (real road-following geometry,
a layered ``casing'' line style, and mode-specific markers), a ranked list of route cards with
fare/duration/walking-distance summaries, a step-by-step segment timeline, and the explanation,
reliability and accessibility panels described in Chapter~5. This is implemented in
\texttt{MultimodalResults.tsx}, \texttt{MapView.tsx}, and \texttt{SegmentTimeline.tsx}.

\section{Dataflow Diagram}

\begin{figure}[H]
  \centering
  \includegraphics[width=0.98\textwidth]{figures/fig-4.4-dfd.png}
  \caption{Data flow diagram (Level 0 / Level 1, combined)}
\end{figure}

\section{Use Case Diagram}

\begin{figure}[H]
  \centering
  \includegraphics[width=0.95\textwidth]{figures/fig-4.1-use-case.png}
  \caption{Use case diagram}
\end{figure}

\section{Class Diagram}

\begin{figure}[H]
  \centering
  \includegraphics[width=\textwidth]{figures/fig-4.2-class.png}
  \caption{Class diagram (core persistence entities)}
\end{figure}

\section{Sequence Diagram}

\begin{figure}[H]
  \centering
  \includegraphics[width=\textwidth]{figures/fig-4.3-sequence.png}
  \caption{Sequence diagram: planning a multimodal route}
\end{figure}

% ============================================================
% CHAPTER 5 — IMPLEMENTATION
% ============================================================
\chapter{Implementation}

\section{Briefing Implementation}

TransitSwap was implemented as a complete, production-shaped full-stack application rather than a
single-script prototype: a React/TypeScript single-page frontend, an Express/TypeScript backend
exposing a 50-endpoint REST API, and MongoDB for optional persistence. The implementation was built
and verified incrementally, phase by phase, with an automated test suite kept green throughout
rather than written retroactively at the end.

\section{System Implementation Phases}

The project was implemented across a sequence of phases, each with its own verification pass:
establishing the baseline frontend scaffold and core routing UI; building the maps-and-routing
foundation on real OSRM/Nominatim data; replacing a fixed four-pattern candidate generator with a
bounded round-based graph search; implementing the accessibility rule engine and its seven
profile-specific hard-constraint rule sets; implementing the weather, crowd and reliability
intelligence layers; implementing the pairwise-logistic-regression personalisation model and fixing
a validity bug that had made its original training labels single-class; building the honest
evaluation harness and sensitivity-analysis tooling used in Chapter~7; hardening performance and
security (response caching, \texttt{helmet}, authenticated evaluation endpoints); completing
user-experience gaps (saved destinations, admin accessibility-report review, explanation
honesty-gating); preparing and executing a live three-part deployment (MongoDB Atlas, Render,
Vercel); and, most recently, fixing a confirmed personalisation bug in which a user's selected
accessibility profile could be silently discarded once enough route-choice history existed to
activate the learned model.

\section{Environmental Support}

Local development runs the frontend via Vite's development server and the backend via
\texttt{ts-node-dev} against either a configured MongoDB Atlas cluster or an ephemeral in-memory
MongoDB instance (\texttt{mongodb-memory-server}) used for both automated tests and local
demonstration without any external database dependency. The deployed system runs the frontend as a
static build on Vercel, the backend as a Node web service on Render, and uses a MongoDB Atlas free
tier cluster for persistence.

\section{Core Implementation Module}

The core of the system is the multimodal routing and ranking pipeline
(\texttt{backend/src/services/multimodalService.ts},
\texttt{accessibilityService.ts}, \texttt{ml/mlPreferenceService.ts},
\texttt{transitDnaService.ts}). Candidate generation combines a small set of fixed journey patterns
with a bounded round-based graph search over the seeded station/stop graph, deduplicates and reduces
the result to its Pareto frontier over (duration, fare, walking distance, transfers), and validates
every candidate's structural and geometric continuity before it is considered further. Accessibility
filtering then removes any candidate that fails the selected profile's hard constraint. Surviving
candidates are enriched with live weather, crowd and reliability data, scored, and explained.

\section{Real-Time Route Intelligence}

Beyond static routing, TransitSwap layers in several real-time and statistically grounded signals:
current weather conditions (OpenWeatherMap, with a clearly labelled demo fallback when unconfigured);
live and recent crowd-level reports blended with a seeded baseline; a data-driven arrival-time
prediction interval computed from real historical duration-prediction errors, with a hard minimum
sample-size gate below which the system reports a rule-of-thumb band and an explicit
\texttt{isDataDriven: false} flag rather than an unearned confidence level; and a 1{,}000-trial
empirical-resampling Monte~Carlo simulation over real historical per-transfer delay samples for
missed-connection risk, rather than an assumed parametric distribution.

\section{Pseudo Code}

\begin{lstlisting}[language=, caption={Candidate generation and ranking (simplified)}]
function generateRoutes(origin, destination, profile):
    candidates <- fixedPatternCandidates(origin, destination)
              + graphSearchCandidates(origin, destination, maxRounds=4)
    candidates <- validateStructuralContinuity(candidates)
    candidates <- removeDuplicates(candidates)
    candidates <- paretoFrontier(candidates, on=[duration, fare, walk, transfers])
    return candidates

function filterAndRank(candidates, profile, userId, preferences):
    accessible <- [ r for r in candidates
                     if not hardConstraintViolated(r, profile) ]
    for r in accessible:
        r.weather      <- getWeatherImpact(r)
        r.crowd        <- getCrowdSummary(r)
        r.reliability  <- getReliabilityScore(r)
        r.confidence   <- getConfidenceInterval(r)   // null if data-driven
                                                      // threshold not met
        r.missedRisk   <- monteCarloMissedConnectionRisk(r, trials=1000)

    weights <- getProfileWeights(profile)
    if hasPersonalizedModel(userId):
        learned <- loadLearnedWeights(userId)
        weights <- floorProfileWeights(learned, weights)   // profile sensitivity
                                                            // can never be diluted
    for r in accessible:
        r.score <- weightedSum(normalize(r.features), weights)

    ranked <- sortDescending(accessible, by=r.score)
    for r in ranked:
        r.explanation <- buildExplanation(r, ranked, profile)  // every claim
                                                                // re-checked here
    return ranked
\end{lstlisting}

% ============================================================
% CHAPTER 6 — TESTING
% ============================================================
\chapter{Testing}

\section{Overview of Testing}

TransitSwap was tested with an extensive automated suite covering the backend domain logic, the
backend's full REST API surface, and the frontend's critical user-facing logic, rather than relying
on manual verification alone. At the time of this report, \textbf{29 backend test files} and
\textbf{35 frontend tests} pass in full, and the suite is re-run after every change to the codebase.

\section{Testing Performed}

\begin{itemize}[leftmargin=*]
  \item \textbf{Unit testing} of individual service functions -- accessibility rule evaluation,
    feature normalisation, weight-floor renormalisation, candidate deduplication, Pareto-frontier
    filtering.
  \item \textbf{Integration testing} of full API endpoints via Supertest against a real (in-memory)
    MongoDB instance -- authentication, route generation, accessibility filtering, crowd and weather
    endpoints, the admin surface, and the research-evaluation endpoint.
  \item \textbf{End-to-end testing} of the complete browser user journey (register, plan a trip,
    save a journey, submit feedback) via Playwright against a real frontend and backend.
  \item \textbf{Regression testing}, re-running the full suite after every feature change described
    in this report, including the most recent accessibility-personalisation fix.
  \item \textbf{Manual exploratory testing} directly in the browser for every accessibility profile,
    used specifically to surface the personalisation bug corrected in the final implementation
    phase.
\end{itemize}

\section{Objective}

The objective of the testing phase was not only to confirm that each function returns a plausible
result, but specifically to confirm that the system's \emph{honesty} properties hold: that an
accessibility-blocked route is actually excluded (not merely down-ranked), that a confidence
interval is never shown without sufficient supporting data, and that an explanation tag is never
displayed unless its underlying claim is verified true for that specific route.

\section{Methodology}

Backend tests are plain Node scripts executed via \texttt{backend/scripts/run-tests.js}, each
asserting a sequence of concrete expectations and printing a pass/fail summary; a database-dependent
test honestly reports itself as \emph{skipped} (not passed) when no MongoDB connection is available,
rather than silently succeeding. Frontend tests use Vitest with React Testing Library. End-to-end
tests spin up a real backend against a freshly created in-memory MongoDB instance and a real Vite
dev server, so the full request path is exercised exactly as a real user would experience it.

\section{Environment}

Tests run locally against Node.js 20, with MongoDB provided either by a configured
\texttt{MONGODB\_URI} or, for fully offline and continuous-integration runs, an ephemeral
\texttt{mongodb-memory-server} instance created fresh for each run.

\section{Result}

All 29 backend test files and all 35 frontend tests pass, and the same checks run automatically on
every push to GitHub (continuous integration). The backend runner counts a database-dependent test that
skipped for lack of a database as a failure in continuous-integration mode, so a passing run means every
database assertion actually executed. Newer regression tests cover the time-of-day matching of crowd data,
the wheelchair step-free rule, and the exclusion reasons returned when no accessible route exists. The
user interface was also scanned with the axe-core accessibility checker against WCAG 2.1 levels A and AA on
every main screen, with no remaining violations.

\begin{table}[H]
\centering
\caption{Testing performed}
\begin{tabular}{@{}p{4cm}p{10cm}@{}}
\toprule
\textbf{Test Type} & \textbf{Description} \\
\midrule
Unit Testing & Accessibility rule evaluation, feature normalisation, weight-floor logic, \\
             & candidate deduplication and Pareto filtering \\
Integration Testing & Full REST API surface (50 endpoints) exercised via Supertest against a \\
                     & real MongoDB instance \\
End-to-End Testing & Complete browser user journey via Playwright (register $\to$ plan $\to$ \\
                    & save $\to$ feedback) \\
Regression Testing & Full suite re-run after every feature change \\
Manual Exploratory Testing & Every accessibility profile tested live in-browser; surfaced the \\
                            & personalisation weight-dilution bug fixed in this project \\
\bottomrule
\end{tabular}
\end{table}

\section{Sample Test Cases}

\begin{table}[H]
\centering
\caption{Sample test cases}
\begin{tabular}{@{}p{1cm}p{5.6cm}p{4.2cm}p{2.6cm}@{}}
\toprule
\textbf{ID} & \textbf{Test Case} & \textbf{Expected Result} & \textbf{Result} \\
\midrule
TC01 & Request a route for a wheelchair profile through a station with a broken lift and uncovered
  stairs & Route is excluded from the result set entirely & Pass \\
TC02 & Submit login with no authentication token & API returns HTTP 401 & Pass \\
TC03 & Request confidence interval with fewer than 5 historical observations & Response reports
  \texttt{isDataDriven: false}, \texttt{confidenceLevelPercent: 0} & Pass \\
TC04 & Select ``stroller'' profile with a route through a known HIGH-crowd station & Response
  includes an explicit HIGH-crowd warning for that profile & Pass \\
TC05 & Personalised user (\(\geq\)5 recorded choices) selects ``pregnant'' profile & Crowd/walking
  weighting is never diluted below the pregnant preset's own value & Pass \\
TC06 & Admin marks a station's lift as broken & Wheelchair-profile route generation immediately
  excludes that station & Pass \\
TC07 & Generate a route with no \texttt{MONGODB\_URI} configured & Route generation still succeeds
  from seeded demo data; auth/saved-journey endpoints return HTTP 503 & Pass \\
TC08 & Reverse-geocode a clicked map coordinate & Returns a real resolved place name and address
  via Nominatim & Pass \\
TC09 & Run the 22-scenario evaluation harness twice with the offline test provider & Produces
  byte-identical results (deterministic) & Pass \\
TC10 & Request route generation with origin equal to destination & Returns an empty result set
  rather than a degenerate zero-length route & Pass \\
\bottomrule
\end{tabular}
\end{table}

% ============================================================
% CHAPTER 7 — RESULTS
% ============================================================
\chapter{Results}

This chapter presents the running application. Screenshots are intentionally left as placeholders
below and should be inserted directly from the live application (local development server or the
deployed instance) before final submission.

\screenshotplaceholder{Landing page}{full-width hero section with the headline search widget}
\screenshotplaceholder{User dashboard}{personalised TransitDNA weights, journey stats, and research-benchmark card}
\screenshotplaceholder{Plan Trip -- accessibility profile selection}{the seven mobility-profile buttons and active-preferences summary}
\screenshotplaceholder{Plan Trip -- map and nearby transit}{interactive Leaflet map with a plotted multimodal route and nearby-station cards}
\screenshotplaceholder{Ranked route candidates with ``Why Recommended'' explanation}{the selected route's explanation tags and the full ranked candidate list}
\screenshotplaceholder{Step-by-step journey timeline}{per-segment walk/bus/metro breakdown with bus route number and boarding-time estimate}
\screenshotplaceholder{Accessibility check and crowd panel}{a route's accessibility warnings (including the HIGH-crowd warning) and station crowd levels}
\screenshotplaceholder{Admin dashboard -- accessibility report review}{pending community accessibility reports and station management}
\screenshotplaceholder{Research / evaluation page}{the 22-scenario benchmark, ablation study, and sensitivity-sweep results}

\bigskip
\noindent In addition to the application screenshots above, the project's own research paper
(produced alongside this report) includes nine further figures generated directly from real
computed data -- pipeline latency by stage, geographic coverage of the seeded network, a real
candidate trade-off scatter for one corridor, the ablation study's decision-change rate per
feature, prediction-interval calibration, a learning curve for the personalisation model, a
sensitivity-analysis tornado chart for the routing cost-model constants, and an example explanation
card -- every one of them traceable to the exact script that produced it.

% ============================================================
% CHAPTER 8 — CONCLUSION
% ============================================================
\chapter{Conclusion}

TransitSwap demonstrates that a multi-criteria, accessibility-first approach to transit routing is
both implementable and meaningfully different in practice from a conventional shortest-time planner
-- not only in principle, but as measured: an ablation study across a 22-scenario benchmark confirms
that every one of six non-time ranking features has a real, measurable effect on the system's final
recommendation, and the system's accessibility engine was verified, live, to immediately change its
output the moment a station's lift status changed. Over the course of this project, a genuine,
user-reported personalisation bug was found and fixed -- a logged-in user's selected accessibility
profile could be silently discarded once enough route-choice history existed for the system's
learned ranking model to activate -- underscoring that a system's stated design intent and its
actual runtime behaviour are not the same thing, and that honest, end-to-end verification (not
assumption) is necessary to close that gap.

The project was also built, throughout, to a standard of transparency beyond simply ``working'':
every reliability and confidence figure the system reports is explicitly labelled as either
genuinely data-driven or a stated rule-of-thumb fallback; every explanation tag shown to a user is
individually checked against the real candidate set before being displayed; and every number
reported in this project's own research-paper evaluation is reproducible from a named script rather
than hand-estimated. The system's real limitations are stated plainly rather than concealed: its
transit network is a synthetic, hand-authored dataset covering a minority of the sampled
metropolitan area; no production-scale user study has been conducted; and several evaluation
figures are necessarily computed on synthetic data because no real-world usage history yet exists.
Future work -- integrating a real GTFS feed, running a genuine user study, and extending calibration
coverage as real usage data accumulates -- is a natural and achievable continuation of the system as
it stands today.

% ============================================================
% REFERENCES
% ============================================================
\chapter*{References}
\addcontentsline{toc}{chapter}{References}
\begin{enumerate}[leftmargin=*]
\item Delling, D., Pajor, T., and Werneck, R. F. (2015). Round-Based Public Transit Routing.
  \textit{INFORMS Journal on Computing}, 27(2), 301--316.
\item Dibbelt, J., Pajor, T., Strasser, B., and Wagner, D. (2018). Connection Scan Algorithm.
  \textit{ACM Journal of Experimental Algorithmics}, 23, 1.7.
\item Burges, C., Shaked, T., Renshaw, E., Lazier, A., Deeds, M., Hamilton, N., and Hullender, G.
  (2005). Learning to Rank using Gradient Descent. \textit{Proceedings of the 22nd International
  Conference on Machine Learning (ICML)}, 89--96.
\item Bast, H., Delling, D., Goldberg, A., Müller-Hannemann, M., Pajor, T., Sanders, P., Wagner,
  D., and Werneck, R. F. (2016). Route Planning in Transportation Networks. In \textit{Algorithm
  Engineering}, Lecture Notes in Computer Science, vol.\ 9220, Springer.
\item Luxen, D., and Vetter, C. (2011). Real-time Routing with OpenStreetMap Data.
  \textit{Proceedings of the 19th ACM SIGSPATIAL International Conference on Advances in
  Geographic Information Systems}, 513--516.
\item Haklay, M., and Weber, P. (2008). OpenStreetMap: User-Generated Street Maps. \textit{IEEE
  Pervasive Computing}, 7(4), 12--18.
\item Koren, Y., Bell, R., and Volinsky, C. (2009). Matrix Factorization Techniques for
  Recommender Systems. \textit{IEEE Computer}, 42(8), 30--37.
\item Guo, Z., and Wilson, N. H. M. (2011). Assessing the Cost of Transfer Inconvenience in Public
  Transport Systems: A Case Study of the London Underground. \textit{Transportation Research
  Part A}, 45(2), 91--104.
\item OpenStreetMap contributors. OpenStreetMap -- \url{https://www.openstreetmap.org}.
\item Nominatim -- OpenStreetMap geocoding service. \url{https://nominatim.org}.
\item Open Source Routing Machine (OSRM). \url{http://project-osrm.org}.
\item React documentation. \url{https://react.dev}.
\item Express.js documentation. \url{https://expressjs.com}.
\item MongoDB / Mongoose documentation. \url{https://www.mongodb.com}, \url{https://mongoosejs.com}.
\item Leaflet -- a JavaScript library for interactive maps. \url{https://leafletjs.com}.
\end{enumerate}

\end{document}
